\chapter{Deploying a Policy on a Robot}
\label{chap:deployment}

Deployment connects a statistical model to a safety-critical dynamical system.
Treat it as an interface and validation problem, not as one command that changes
\filepath{device=cpu} to \filepath{device=cuda}.

\section{Implement a robot adapter}

For a new robot, isolate four responsibilities:

\begin{enumerate}
  \item build observations in the exact training format;
  \item load normalization and run the policy;
  \item decode predicted actions into Cartesian or joint targets;
  \item enforce low-level control limits and stop conditions.
\end{enumerate}

The policy should not directly open sockets, read cameras, or implement robot
kinematics. A robot-specific adapter makes it possible to reuse the same policy
runner and replace only state acquisition and command execution.

\begin{projectbox}[UR10 implementation]
\filepath{ur10_real_robot/policy_exec/} contains the observation builder and
real runner. The backend reads UR10 state, two RealSense cameras, and gripper
width; a Pinocchio Cartesian servo converts pose targets to bounded joint
velocities; the gripper uses a Modbus adapter.
\end{projectbox}

\section{Deploy in stages}

Never begin with unrestricted physical motion. Use at least these stages:

\begin{description}
  \item[1. Offline/fake:] load the checkpoint, create correctly shaped dummy
  observations, and inspect predicted actions.
  \item[2. Real observations, no motion:] connect cameras and state readers,
  run inference, and print targets while commands remain disabled.
  \item[3. Low-speed physical motion:] use conservative limits, one object pose,
  a clear workspace, and an operator ready to stop.
  \item[4. Frozen evaluation:] fix checkpoint and parameters, then execute the
  written protocol.
\end{description}

At each stage, compare current observations with training data and verify units,
frames, quaternion order, TCP offset, action horizon, and gripper convention.

\section{Synchronous and asynchronous inference}

If inference is slower than action execution, synchronous replanning causes the
robot to pause while each new sequence is predicted. Asynchronous execution can
hide this delay:

\begin{figure}[H]
  \centering
  \includegraphics[width=0.96\textwidth]{asynchronous_inference_timeline.png}
  \caption{Asynchronous receding-horizon execution. The current buffered plan
  continues while the next plan is computed, then execution switches when the
  new plan is ready.}
  \label{fig:async-manual}
\end{figure}

This improves continuity but does not remove latency. The new plan is based on
an observation captured before inference began, so it may already be stale when
execution starts. Faster robot motion, longer inference, and contact-rich tasks
increase this mismatch.

Measure:

\begin{itemize}
  \item policy inference latency and its variance;
  \item age of the observation used for each plan;
  \item remaining actions when a new plan becomes ready;
  \item target discontinuity at plan switches;
  \item camera and control-loop delays separately.
\end{itemize}

Possible improvements include a deployment accelerator, a smaller model, a
faster valid decoder or sampler, temporal ensembling, latency-aware action
selection, and blending compatible plan transitions. Revalidate behavior after
changing model-specific inference settings.

\section{Parameters have different responsibilities}

Do not tune every parameter as a generic ``speed'' control.

\begin{table}[H]
  \centering
  \caption{Main deployment parameter groups.}
  \begin{tabularx}{\textwidth}{@{}p{0.32\textwidth}X@{}}
    \toprule
    \textbf{Group} & \textbf{What it controls} \\
    \midrule
    Model decoding/sampling & Computation cost and output approximation. \\
    Action/execution horizon & How long one prediction is consumed before replacement. \\
    Action execution rate & Time spacing between buffered targets. \\
    Cartesian gains & How aggressively tracking error is corrected. \\
    Target smoothing & Discontinuity/noise versus response lag. \\
    Joint and Cartesian limits & Maximum physical motion allowed by software. \\
    Position/rotation stop thresholds & Maximum discrepancy accepted before stopping. \\
    Teach-pendant slider & Hardware-level scaling of physical speed. \\
    \bottomrule
  \end{tabularx}
\end{table}

Record both requested software limits and the teach-pendant slider. Two runs
with the same command but different slider values are not the same execution
condition.

\section{Safety stops should latch and explain themselves}

A stop should cancel motion, prevent the same invalid target from being sent
repeatedly, and print enough context to diagnose the cause. Useful checks include:

\begin{itemize}
  \item target position and orientation error;
  \item joint velocity and workspace limits;
  \item stale camera or robot state;
  \item excessive control-loop time;
  \item loss of robot or gripper communication;
  \item invalid or non-finite policy output.
\end{itemize}

Do not increase a threshold only because it interrupts a trial. First determine
whether the large error comes from a wrong frame, quaternion, delayed plan,
unreachable target, or genuine tracking lag.

\section{Camera and gripper behavior during deployment}

Use the same camera order, crop, resolution, and exposure strategy as training.
Warm up cameras before the first observation and reject stale frames. If the
scene requires auto-exposure, give the camera time to adapt in both
demonstrations and deployment.

Decode the gripper action using the same semantics used during recording. Log
raw policy output, post-processing, command sent, target width, and measured
width. This distinguishes policy hesitation from actuator delay or contact.

\section{Use a staged execution wrapper}

The repository wrapper \filepath{ur10_real_robot/run_diffusion_real.sh}
defaults to motion disabled and requires an explicit motion-enable confirmation.
Its validated project commands are in \filepath{README_CODE.md}. A new robot
adapter should expose equivalent fake, observation-only, and physical-motion
stages.

\begin{safetybox}[Learned policy execution]
A successful checkpoint load does not imply a safe target. Keep independent
software limits, robot safety configuration, low initial speed, an unobstructed
workspace, and an operator able to stop motion immediately.
\end{safetybox}

\begin{checkpointbox}[Ready for evaluation]
The selected checkpoint and command are frozen; camera images match the dataset;
timing has been measured; repeated low-speed trials show bounded tracking; and
every stop condition has been tested or reviewed.
\end{checkpointbox}
